\begin{proof}
Follows from Property~\ref{prop2.5_C5} in Proposition~\ref{prop:dc_color}.
\end{proof}

\setbox\toto=\hbox{See~\cite[Section~3]{GK}}
\begin{lemm}[\box\toto]
\label{lem:fc}
%\textup{(See~\cite[Section~3]{GK})}
Let $v \in W$ be a fully commutative element 
and $v = s_{i_1} \cdots s_{i_r} = s_{j_1} \cdots s_{j_r}$ be its reduced expressions.
Then we have
\begin{enumerate}[(a)]
\item\label{lemma3.6_a} %[(a)]
Let $i$, $j$ be adjacent nodes in the Dynkin diagram.
Then the subsequence of $(i_1, \ldots, i_r)$ consisting of $i$ and $j$ 
is identical with the subsequence of $(j_1, \ldots, j_r)$ consisting of $i$ and $j$.
\item\label{lemma3.6_b} %[(b)]
Let $i$ be a node in the Dynkin diagram.
Then the number of occurrence of $i$ in $(i_1, \ldots, i_r)$ 
is equal to the number of occurrence of $i$ in $(j_1, \ldots, j_r)$.
\end{enumerate}
\end{lemm}

We use these lemmas to prove the characterization of excited diagrams.

\begin{proof}[Proof of Proposition~\ref{prop:E=R}]
We follow the same idea used in the proof of~\cite[Proposition~4.8]{GK}.
We denote by $\RR_P(F)$ the set of all subsets $D \subset P$ satisfying $\# D = \# F$ and $w_D = w_F$.

First we prove that $\EE_P(F) \subset \RR_P(F)$.
Since $F \in \RR_P(F)$, it is enough to show that, 
if $D' \in \EE_P(F)$ is obtained from $D \in \EE_P(F)$ by an elementary excitation, 
then $w_{D'} = w_D$.
Let $u$, $v \in P$ be elements such that $[v,u]$ is a $d_k$-interval, 
$[v,u] \cap D \cap N_{c(u)} = \emptyset$ and 
$D' = \alpha_u(D) = D \setminus \{ u \} \cup \{ v \}$.
If $v = p_k$, $u = p_l$ and $\{ j : k < j < l, \ p_j \in D \} = \{ j_1, \dots, j_m \}$ 
($j_1 < \dots < j_m$), then we have
$$
w_D = \cdots s(p_{j_1}) \cdots s(p_{j_m}) s(p_l) \cdots,
\quad
w_{D'} = \cdots s(p_k) s(p_{j_1}) \cdots s(p_{j_m}) \cdots.
$$
Then by using Lemma~\ref{lem:comm}, we have $w_{D'} = w_D$.

Next we prove that $\RR_P(F) \subset \EE_P(F)$.
Since $F$ is an order filter, we can take a linear extension of $P$ 
such that $F = \{ p_{n+1}, \ldots, p_N \}$, where $n = \# (P \setminus F)$.
We define the energy $e(D)$ of any subset $D \subset P$ with $\# D = \# F$ by putting
$$
e(D) = \sum_{p_k \in F} k - \sum_{p_k \in D} k.
$$
Then, by the assumption on our linear extension, 
we see that $e(D) \ge 0$ and that $e(D) = 0$ if and only if $D = F$.

We proceed by induction on $e(D)$ to prove $D \in \RR_P(F)$ implies $D \in \EE_P(F)$.
Let $D \in \RR_P(F)$.
If $e(D) = 0$, then we have $D = F \in \EE_P(F)$.

Suppose that $e(D) > 0$.
Then there exists an element $u \in F$ such that $u \not\in D$.
Let $u$ be the last element, i.e. $u = p_l$ with $l$ largest, 
satisfying $u \in F$ and $u \not\in D$.
Since $w_D = w_F$ and $w_F$ is fully commutative, it follows from Lemma~\ref{lem:fc}~\ref{lemma3.6_b} %
that there exists an element $v \in D$ with the same color $i$ as $u$.
Let $v = p_k$ ($1 \le k < l$) be the last element satisfying $v \in D$ and $c(v) = c(u) = i$.
Then, by the choice of $k$ and $l$, we see that the reduced expressions of $w_F$ and $w_D$
are of the form $w_F = \cdots s_i s_{h_1} \cdots s_{h_r}$ ($s_i$ corresponds to $u$) 
and $w_D = \cdots s_i s_{j_1} \cdots s_{j_m} s_{h_1} \cdots s_{h_r}$ ($s_i$ corresponds to $v$) 
with no $i$ appearing in the segment $(j_1, \dots, j_m)$.
If $z \in [v,u] \cap D \cap N_i$, then $j = c(z)$ appears in the segment $(j_1, \dots, j_m)$ 
and this contradicts to the assertion of Lemma~\ref{lem:fc}~\ref{lemma3.6_a}.
Hence we have $[v,u] \cap D \cap N_i = \emptyset$.
If we put $D' = D \setminus \{ v \} \cup \{ u \}$,
then $w_{D'} = w_F$ by Lemma~\ref{lem:comm} and $e(D') < e(D)$.
Then by the induction hypothesis we have $D' \in \EE_P(F)$.
Since $D$ is obtained from $D'$ by a sequence of elementary excitations, 
we obtain $D \in \EE_P(F)$.
\end{proof}

Next we give a non-recursive description of the set of excited peaks $B(D)$, 
which implies that $B(D)$ is well-defined, i.e. 
it is independent of the choices of elementary excitations to reach $D$ from $F$.

\begin{prop}
\label{prop:BD}
Let $D \in \EE_P(F)$ be an excited diagram.
\begin{enumerate}[(a)]
\item\label{prop3.7_a} %[(a)]
The following are equivalent for $x \in P$:
\begin{enumerate}[(i)]
\item\label{prop3.7_a_i} %[(i)]
$x \in B(D)$.
\item\label{prop3.7_a_ii} %[(ii)]
There exists an element $y \in D_{c(x)}$ such that $y<x$ and $[y,x] \cap D \cap N_{c(x)} = \emptyset$.
\end{enumerate}
\item\label{prop3.7_b} %[(b)]
We have $D \cap B(D) = \emptyset$.
\end{enumerate}
\end{prop}

In the proof of this proposition, we utilize the following lemma, 
which will be used also in the sequel of this section.

\begin{lemm}
\label{lem:interval}
Let $x$, $y$, $z$, $u$ and $v$ be elements of $P$ such that $c(x) = c(y) = c(z) = i$, 
$c(u) = c(v) = j$,
$z<y<x$ and $v<u$.
Let $D$ be a subset of $P$.
Then we have
\begin{enumerate}[(a)]
\item\label{lemma3.8_a} %[(a)]
If $[z,y] \cap D \cap N_i = \emptyset$ and $[y,x] \cap D \cap N_i = \emptyset$, 
then we have $[z,x] \cap D \cap N_i = \emptyset$.
\item\label{lemma3.8_b} %[(b)]
Suppose $u \in D$ and $v \not\in D$.
If $[y,x] \cap D \cap N_i = \emptyset$ and $x \not\in [v,u] \cap N_j$, 
then we have $[y,x] \cap (D \cup \{ v \}) \cap N_i = \emptyset$.
\item\label{lemma3.8_c} %[(c)]
Suppose $u \in D$ and $v \not\in D$.
If $[y,x] \cap (D \setminus \{ u \} \cup \{ v \}) \cap N_i = \emptyset$ 
and $y \not\in [v,u] \cap N_j$, 
then we have $[y,x] \cap D \cap N_i = \emptyset$.
\item\label{lemma3.8_d} %[(d)]
Suppose that $u \not\in D$ and $v \in D$.
If $[y,x] \cap D \cap N_i = \emptyset$ and $y \not\in [v,u] \cap N_j$, 
then we have $[y,x] \cap (D \cup \{ u \}) \cap N_i = \emptyset$.
\end{enumerate}
\end{lemm}

\begin{proof}
%%(a)
\ref{lemma3.8_a}~By using Property~\ref{prop2.5_C5} in Proposition~\ref{prop:dc_color}, 
we have $[z,x] \cap N_i = ( [z,y] \cap N_i ) \cup ( [y,x] \cap N_i)$.

%(b)
\ref{lemma3.8_b}~Assume to the contrary that $[y,x] \cap (D \cup \{ v \}) \cap N_i \neq \emptyset$.
Since $[y,x] \cap D \cap N_i = \emptyset$, we have $v \in [y,x] \cap N_i$.
Thus $j = c(u) = c(v)$ is adjacent to $i = c(x)$ in $\Gamma$.
Since $u \in D$ and $[y,x] \cap D \cap N_i = \emptyset$, we have $u \not\in [y,x] \cap N_i$.
Hence, by using Property~\ref{prop2.5_C5} in Proposition~\ref{prop:dc_color},
we see that $y < v < x < u$ and $x \in [v,u] \cap N_j$, 
which contradicts to the assumption.

%(c)
\ref{lemma3.8_c}~By an argument similar to~\ref{lemma3.8_b}, we can show that, if $[y,x] \cap D \cap N_i \neq \emptyset$, 
then $v < y < u < x$, which contradicts to $y \not\in [v,u] \cap N_j$.

%(d)
\ref{lemma3.8_d}~By an argument similar to~\ref{lemma3.8_b}, we can show that, if $[y,x] \cap (D \cup \{ u \}) \cap N_i \neq \emptyset$, 
then $v < y < u < x$, which contradicts to $y \not\in [v,u] \cap N_j$.
\end{proof}

\begin{proof}[Proof of Proposition~\ref{prop:BD}]
We denote by $B'(D)$ the subset of $P$ consisting of elements $x \in P$ satisfying the condition~\ref{prop3.7_a_ii} in~\ref{prop3.7_a},
and prove $B(D) = B'(D)$ and $D \cap B'(D) = \emptyset$.
We proceed by induction on the number of elementary excitations to reach $D$ from~$F$.

We begin with considering the case where $D = F$.
Let $x$ and $y$ be elements of $F$ with the same color $i$ satisfying $y<x$.
Then it follows from Properties~\ref{prop2.5_C4} and~\ref{prop2.5_C5} in Proposition~\ref{prop:dc_color} that 
an element $z$ covered by $x$ or covers $y$ belongs to $[y,x] \cap F \cap N_i$.
Hence we have $B'(F) = \emptyset = B(F)$ and $F \cap B'(F) = \emptyset$.

We prove that $B(\alpha_u(D)) = B'(\alpha_u(D))$ and $\alpha_u(D) \cap B'(\alpha_u(D)) = \emptyset$ 
for $D \in \EE_P(F)$ and a $D$-active element $u \in D$.
Let $v$ be the element such that $v<u$,
$c(v)=c(u)$, $[v,u] \cap D \cap N_{c(u)} = \emptyset$,
$[v,u]$ is a $d_k$-interval
and $\alpha_u(D) = D \setminus \{ u \} \cup \{ v \}$.
Recall that, by definition, 
$B(\alpha_u(D)) = \left( B(D) \setminus ([v,u] \cap N_{c(u)}) \right) \cup \{ u \}$.

First we show $B(\alpha_u(D)) \subset B'(\alpha_u(D))$.
Since $v \in \alpha_u(D)$ and $[v,u] \cap D \cap N_{c(u)} = \emptyset$, 
we have $u \in B'(\alpha_u(D))$.
Let $x \in B(D)$ such that $x \not\in [v,u] \cap N_{c(u)}$.
Since $B(D) = B'(D)$ by the induction hypothesis, 
there exists $y \in D_{c(x)}$ such that $y<x$ and $[y,x] \cap D \cap N_{c(x)} = \emptyset$.
Then, by using Lemma~\ref{lem:interval}~\ref{lemma3.8_b}, 
we have $[y,x] \cap \alpha_u(D) \cap N_{c(x)} = \emptyset$, hence $x \in B'(\alpha_u(D))$.

Next, in order to show $B'(\alpha_u(D)) \subset B(\alpha_u(D))$, 
we take an element $x \in B'(\alpha_u(D))$ such that $x \neq u$ 
and prove $x \in B(D) \setminus ([v,u] \cap N_{c(u)})$.
Then there exists $y \in (\alpha_u(D))_{c(x)}$ such that 
$y<x$ and $[y,x] \cap \alpha_u(D) \cap N_{c(x)} = \emptyset$.
If $y = v$, then $\alpha_u(D)\cap N_{c(x)}=D\cap N_{c(x)}$, 
hence $[u,x] \cap D \cap N_{c(x)} \subset [v,x] \cap \alpha_u(D) \cap N_{c(x)} = \emptyset$.
Since $u \in D$, we have $x \in B'(D)=B(D)$ and $x \not\in N_{c(u)}$.
We consider the case where $y \neq v$.
In this case, $y \in D, y\not\in [v,u]\cap N_{c(u)}$ 
and it follows from Lemma~\ref{lem:interval}~\ref{lemma3.8_c} 
that $[y,x] \cap D \cap N_{c(x)} = \emptyset$,
thus $x\in B'(D)=B(D)$.
Also we have $x \not\in [v,u] \cap N_{c(u)}$.
In fact, if $x \in [v,u] \cap N_{c(u)}$, 
then $c(y) = c(x)$ is adjacent to $c(u)$ 
and it follows from $y \in D$ and $[v,u] \cap D \cap N_{c(u)} = \emptyset$ that $y \not\in [v,u]\cap N_{c(u)}$.
Hence, by using Property~\ref{prop2.5_C5} in Proposition~\ref{prop:dc_color}, 
we have $y < v < x < u$ and 
this contradicts to $[y,x] \cap \alpha_u(D) \cap N_{c(x)} = \emptyset$.
Therefore we have $x \in B(D) \setminus ([v,u] \cap N_{c(u)})$.

Finally we show that $\alpha_u(D) \cap B'(\alpha_u(D)) = \emptyset$.
Let $x$ and $y$ be elements of $\alpha_u(D)$ with the same color $i$ satisfying $y<x$.
Since $\alpha_u(D) = D \setminus \{ u \} \cup \{ v\}$, it is enough to 
show 
$[y,x] \cap \alpha_u(D) \cap N_i \neq \emptyset$
for the following three cases:
$$
\text{Case~1. $x$, $y \in D$,}
\quad
\text{Case~2. $y = v$,}
\quad
\text{Case~3. $x = v$.}
$$
In Case~1, assume that $[y,x] \cap \alpha_u(D) \cap N_i = \emptyset$.
If $[y,x] \cap D \cap N_i = \emptyset$, then we have $x \in B(D)$ 
and this contradicts to the induction hypothesis $D \cap B(D) = \emptyset$.
Hence $[y,x] \cap D \cap N_i \neq \emptyset$ 
and this implies $u \in [y,x] \cap N_i$, 
i.e. $y < u < x$ and $c(u) = c(v)$ is adjacent to $i$ in $\Gamma$.
Since $v \in \alpha_u(D)$, we have $v \not\in [y,x] \cap N_i$.
Therefore, by using Property~\ref{prop2.5_C5} in Proposition~\ref{prop:dc_color}, we see that 
$v < y < u < x$, which contradicts to the $D$-activity of $u$.
In Case~2, we have $v < u < x$ because $[v,u]$ is a $d_k$-interval 
and there is no element $w \in [v,u]$ with $w \neq v, u$ and $c(w) = c(v) = c(u)$ 
by Property~\ref{prop2.5_C2} in Proposition~\ref{prop:dc_color}.
Since $[u,x] \cap D \cap N_i \neq \emptyset$ by the induction hypothesis, 
we have $[v,x] \cap \alpha_u(D) \cap N_i \neq \emptyset$.
In Case~3, we have $[v,u] \cap D \cap N_i = \emptyset$ ($u$ is $D$-active) 
and $[y,u] \cap D \cap N_i \neq \emptyset$ (the induction hypothesis).
Hence by using Lemma~\ref{lem:interval}~\ref{lemma3.8_a} we obtain $[y,v] \cap \alpha_u(D) \cap N_i \neq \emptyset$.
Therefore we see that any element satisfying the condition~\ref{prop3.7_a_ii} in~\ref{prop3.7_a} for $\alpha_u(D)$ 
does not belong to $\alpha_u(D)$, and $\alpha_u(D) \cap B'(\alpha_u(D)) = \emptyset$.
This completes the proof.
\end{proof}

\subsection{%
\texorpdfstring{$K$}{K}-theoretical excited diagrams
}

We define $K$-theoretical excited diagrams and study their properties.
For shapes and shifted shapes, these diagrams were introduced in~\cite{GK}.

\begin{defi}
\label{def:K-excited}
Let $P$ be a $d$-complete poset and let $F$ be an order filter of $P$.
\begin{enumerate}[(a)]
\item~\label{defi3.9_a} %[(a)]
Let $D$ be a subset of $P$ and $u \in D$.
If $u$ is $D$-active and $[v,u]$ is a $d_k$-interval, 
then we define $\alpha^*_u(D)$ to be the subset of $P$ obtained by 
adding $v$ to $D$.
We call this operation a \emph{$K$-theoretical elementary excitation}.
\item~\label{defi3.9_b} %[(b)]
A \emph{$K$-theoretical excited diagram} of $F$ in $P$ is a subset of $P$ 
obtained from $F$ after a sequence of ordinary and $K$-theoretical elementary excitations on active elements. 
Let $\EE^*_P(F)$ be the set of all $K$-theoretical excited diagrams of $F$ in $P$.
\end{enumerate}
\end{defi}

\begin{exam}
\label{ex:shifted_K-excited}
If $P = S(5,4,2,1)$ is the shifted shape corresponding to a strict partition $(5,4,2,1)$ 
and $F = S(3,1)$, then there are $11$ $K$-theoretical excited diagrams in $\EE^{*}_P(F)$ 
shown in Figure~\ref{fig:shifted_K-excited}, 
and five of them are ordinary excited diagrams in $\EE_P(F)$.
In Figure~\ref{fig:shifted_K-excited}, 
the shaded cells form a $K$-theoretical exited diagram, 
and the arrow $D \longrightarrow D'$ (\resp $D \Longrightarrow D'$) indicates that 
$D'$ is obtained from $D$ by an ordinary (\resp $K$-theoretical) elementary excitation.
\begin{figure}[!htb]
$$
\setlength{\unitlength}{1.1pt}
\xymatrix@!C=9mm{
&&&
\raisebox{-20pt}{
\begin{picture}(50,40)
\fboxsep=0mm
\put(0,30){\colorbox[gray]{0.7}{\makebox(10,10){}}}
\put(10,30){\colorbox[gray]{0.7}{\makebox(10,10){}}}
\put(20,30){\colorbox[gray]{0.7}{\makebox(10,10){}}}
\put(10,20){\colorbox[gray]{0.7}{\makebox(10,10){}}}
\put(0,40){\line(1,0){50}}
\put(0,30){\line(1,0){50}}
\put(10,20){\line(1,0){40}}
\put(20,10){\line(1,0){20}}
\put(30,0){\line(1,0){10}}
\put(0,30){\line(0,1){10}}
\put(10,20){\line(0,1){20}}
\put(20,10){\line(0,1){30}}
\put(30,0){\line(0,1){40}}
\put(40,0){\line(0,1){40}}
\put(50,20){\line(0,1){20}}
\end{picture}
}
\ar[llld]
\ar[ld]
\ar@{=>}[rd]
\ar@{=>}[rrrd]
\\
\raisebox{-20pt}{
\begin{picture}(50,40)
\fboxsep=0mm
\put(0,30){\colorbox[gray]{0.7}{\makebox(10,10){}}}
\put(10,30){\colorbox[gray]{0.7}{\makebox(10,10){}}}
\put(30,20){\colorbox[gray]{0.7}{\makebox(10,10){}}}
\put(10,20){\colorbox[gray]{0.7}{\makebox(10,10){}}}
\put(0,40){\line(1,0){50}}
\put(0,30){\line(1,0){50}}
\put(10,20){\line(1,0){40}}
\put(20,10){\line(1,0){20}}
\put(30,0){\line(1,0){10}}
\put(0,30){\line(0,1){10}}
\put(10,20){\line(0,1){20}}
\put(20,10){\line(0,1){30}}
\put(30,0){\line(0,1){40}}
\put(40,0){\line(0,1){40}}
\put(50,20){\line(0,1){20}}
\end{picture}
}
\ar[d]
\ar@{=>}[rrd]
&&
\raisebox{-20pt}{
\begin{picture}(50,40)
\fboxsep=0mm
\put(0,30){\colorbox[gray]{0.7}{\makebox(10,10){}}}
\put(10,30){\colorbox[gray]{0.7}{\makebox(10,10){}}}
\put(20,30){\colorbox[gray]{0.7}{\makebox(10,10){}}}
\put(30,0){\colorbox[gray]{0.7}{\makebox(10,10){}}}
\put(0,40){\line(1,0){50}}
\put(0,30){\line(1,0){50}}
\put(10,20){\line(1,0){40}}
\put(20,10){\line(1,0){20}}
\put(30,0){\line(1,0){10}}
\put(0,30){\line(0,1){10}}
\put(10,20){\line(0,1){20}}
\put(20,10){\line(0,1){30}}
\put(30,0){\line(0,1){40}}
\put(40,0){\line(0,1){40}}
\put(50,20){\line(0,1){20}}
\end{picture}
}
\ar[lld]
\ar@{=>}[rrd]
&&
\raisebox{-20pt}{
\begin{picture}(50,40)
\fboxsep=0mm
\put(0,30){\colorbox[gray]{0.7}{\makebox(10,10){}}}
\put(10,30){\colorbox[gray]{0.7}{\makebox(10,10){}}}
\put(20,30){\colorbox[gray]{0.7}{\makebox(10,10){}}}
\put(30,20){\colorbox[gray]{0.7}{\makebox(10,10){}}}
\put(10,20){\colorbox[gray]{0.7}{\makebox(10,10){}}}
\put(0,40){\line(1,0){50}}
\put(0,30){\line(1,0){50}}
\put(10,20){\line(1,0){40}}
\put(20,10){\line(1,0){20}}
\put(30,0){\line(1,0){10}}
\put(0,30){\line(0,1){10}}
\put(10,20){\line(0,1){20}}
\put(20,10){\line(0,1){30}}
\put(30,0){\line(0,1){40}}
\put(40,0){\line(0,1){40}}
\put(50,20){\line(0,1){20}}
\end{picture}
}
\ar[d]
\ar@{=>}[rrd]
&&
\raisebox{-20pt}{
\begin{picture}(50,40)
\fboxsep=0mm
\put(0,30){\colorbox[gray]{0.7}{\makebox(10,10){}}}
\put(10,30){\colorbox[gray]{0.7}{\makebox(10,10){}}}
\put(20,30){\colorbox[gray]{0.7}{\makebox(10,10){}}}
\put(10,20){\colorbox[gray]{0.7}{\makebox(10,10){}}}
\put(30,00){\colorbox[gray]{0.7}{\makebox(10,10){}}}
\put(0,40){\line(1,0){50}}
\put(0,30){\line(1,0){50}}
\put(10,20){\line(1,0){40}}
\put(20,10){\line(1,0){20}}
\put(30,0){\line(1,0){10}}
\put(0,30){\line(0,1){10}}
\put(10,20){\line(0,1){20}}
\put(20,10){\line(0,1){30}}
\put(30,0){\line(0,1){40}}
\put(40,0){\line(0,1){40}}
\put(50,20){\line(0,1){20}}
\end{picture}
}
\ar[lllld]
\ar@{=>}[d]
\\
\raisebox{-20pt}{
\begin{picture}(50,40)
\fboxsep=0mm
\put(0,30){\colorbox[gray]{0.7}{\makebox(10,10){}}}
\put(10,30){\colorbox[gray]{0.7}{\makebox(10,10){}}}
\put(30,20){\colorbox[gray]{0.7}{\makebox(10,10){}}}
\put(30,00){\colorbox[gray]{0.7}{\makebox(10,10){}}}
\put(0,40){\line(1,0){50}}
\put(0,30){\line(1,0){50}}
\put(10,20){\line(1,0){40}}
\put(20,10){\line(1,0){20}}
\put(30,0){\line(1,0){10}}
\put(0,30){\line(0,1){10}}
\put(10,20){\line(0,1){20}}
\put(20,10){\line(0,1){30}}
\put(30,0){\line(0,1){40}}
\put(40,0){\line(0,1){40}}
\put(50,20){\line(0,1){20}}
\end{picture}
}
\ar[d]
\ar@{=>}[rrd]
&&
\raisebox{-20pt}{
\begin{picture}(50,40)
\fboxsep=0mm
\put(0,30){\colorbox[gray]{0.7}{\makebox(10,10){}}}
\put(10,30){\colorbox[gray]{0.7}{\makebox(10,10){}}}
\put(10,20){\colorbox[gray]{0.7}{\makebox(10,10){}}}
\put(30,20){\colorbox[gray]{0.7}{\makebox(10,10){}}}
\put(30,0){\colorbox[gray]{0.7}{\makebox(10,10){}}}
\put(0,40){\line(1,0){50}}
\put(0,30){\line(1,0){50}}
\put(10,20){\line(1,0){40}}
\put(20,10){\line(1,0){20}}
\put(30,0){\line(1,0){10}}
\put(0,30){\line(0,1){10}}
\put(10,20){\line(0,1){20}}
\put(20,10){\line(0,1){30}}
\put(30,0){\line(0,1){40}}
\put(40,0){\line(0,1){40}}
\put(50,20){\line(0,1){20}}
\end{picture}
}
&&
\raisebox{-20pt}{
\begin{picture}(50,40)
\fboxsep=0mm
\put(0,30){\colorbox[gray]{0.7}{\makebox(10,10){}}}
\put(10,30){\colorbox[gray]{0.7}{\makebox(10,10){}}}
\put(20,30){\colorbox[gray]{0.7}{\makebox(10,10){}}}
\put(30,20){\colorbox[gray]{0.7}{\makebox(10,10){}}}
\put(30,0){\colorbox[gray]{0.7}{\makebox(10,10){}}}
\put(0,40){\line(1,0){50}}
\put(0,30){\line(1,0){50}}
\put(10,20){\line(1,0){40}}
\put(20,10){\line(1,0){20}}
\put(30,0){\line(1,0){10}}
\put(0,30){\line(0,1){10}}
\put(10,20){\line(0,1){20}}
\put(20,10){\line(0,1){30}}
\put(30,0){\line(0,1){40}}
\put(40,0){\line(0,1){40}}
\put(50,20){\line(0,1){20}}
\end{picture}
}
&&
\raisebox{-20pt}{
\begin{picture}(50,40)
\fboxsep=0mm
\put(0,30){\colorbox[gray]{0.7}{\makebox(10,10){}}}
\put(10,30){\colorbox[gray]{0.7}{\makebox(10,10){}}}
\put(30,20){\colorbox[gray]{0.7}{\makebox(10,10){}}}
\put(10,20){\colorbox[gray]{0.7}{\makebox(10,10){}}}
\put(20,30){\colorbox[gray]{0.7}{\makebox(10,10){}}}
\put(30,0){\colorbox[gray]{0.7}{\makebox(10,10){}}}
\put(0,40){\line(1,0){50}}
\put(0,30){\line(1,0){50}}
\put(10,20){\line(1,0){40}}
\put(20,10){\line(1,0){20}}
\put(30,0){\line(1,0){10}}
\put(0,30){\line(0,1){10}}
\put(10,20){\line(0,1){20}}
\put(20,10){\line(0,1){30}}
\put(30,0){\line(0,1){40}}
\put(40,0){\line(0,1){40}}
\put(50,20){\line(0,1){20}}
\end{picture}
}
\\
\raisebox{-20pt}{
\begin{picture}(50,40)
\fboxsep=0mm
\put(0,30){\colorbox[gray]{0.7}{\makebox(10,10){}}}
\put(20,20){\colorbox[gray]{0.7}{\makebox(10,10){}}}
\put(30,20){\colorbox[gray]{0.7}{\makebox(10,10){}}}
\put(30,0){\colorbox[gray]{0.7}{\makebox(10,10){}}}
\put(0,40){\line(1,0){50}}
\put(0,30){\line(1,0){50}}
\put(10,20){\line(1,0){40}}
\put(20,10){\line(1,0){20}}
\put(30,0){\line(1,0){10}}
\put(0,30){\line(0,1){10}}
\put(10,20){\line(0,1){20}}
\put(20,10){\line(0,1){30}}
\put(30,0){\line(0,1){40}}
\put(40,0){\line(0,1){40}}
\put(50,20){\line(0,1){20}}
\end{picture}
}
&&
\raisebox{-20pt}{
\begin{picture}(50,40)
\fboxsep=0mm
\put(0,30){\colorbox[gray]{0.7}{\makebox(10,10){}}}
\put(10,30){\colorbox[gray]{0.7}{\makebox(10,10){}}}
\put(20,20){\colorbox[gray]{0.7}{\makebox(10,10){}}}
\put(30,20){\colorbox[gray]{0.7}{\makebox(10,10){}}}
\put(30,0){\colorbox[gray]{0.7}{\makebox(10,10){}}}
\put(0,40){\line(1,0){50}}
\put(0,30){\line(1,0){50}}
\put(10,20){\line(1,0){40}}
\put(20,10){\line(1,0){20}}
\put(30,0){\line(1,0){10}}
\put(0,30){\line(0,1){10}}
\put(10,20){\line(0,1){20}}
\put(20,10){\line(0,1){30}}
\put(30,0){\line(0,1){40}}
\put(40,0){\line(0,1){40}}
\put(50,20){\line(0,1){20}}
\end{picture}
}
}
$$
\caption{$K$-theoretical excited diagrams of $S(3,1)$ in $S(5,4,2,1)$}
\label{fig:shifted_K-excited}
\end{figure}
\end{exam}


For a fixed linear extension of $P$ and 
a subset $D = \{ p_{i_1}, \ldots, p_{i_r} \}$ ($i_1 < \cdots < i_r$) of $P$, 
we define an element $w^*_D \in W$ by putting
$$
w^*_D = s(p_{i_1}) * s(p_{i_2}) * \cdots * s(p_{i_r}),
$$
where $* : W \times W \to W$ is the associative product, called the Demazure product, defined by
$$
s_i * w
 =
\begin{cases}
 s_i w &\text{if $l(s_i w) = l(w) + 1$,} \\
 w &\text{if $l(s_i w) = l(w) - 1$.}
\end{cases}
$$
For the fundamental properties of the Demazure product
(also called the Hecke product), we refer to~\cite[Section 3]{BM}.
Since $w_P$ is fully commutative (Proposition~\ref{prop:w_P}), the element $w^*_D$ is 
independent of the choices of linear extensions of $P$.

The following proposition is a key to rephrase the Billey-type formula for equivariant $K$-theory 
in terms of combinatorics of $d$-complete posets (see Proposition~\ref{prop:Billey_P}).

\begin{prop}
\label{prop:E*=R*}
Let $P$ be a connected $d$-complete poset and $F$ an order filter of $P$.
Then a subset $D \subset P$ is a $K$-theoretical excited diagram of $F$ in $P$ 
if and only if $w^*_D = w_F$.
\end{prop}

We follow the same idea as the proof of~\cite[Proposition~4.8]{GK}.

\setbox\toto=\hbox{\cite[Lemma~3.1 and Proposition~3.4]{GK}}
\begin{lemm}[\box\toto]
\label{lem:red}
%\textup{(\cite[Lemma~3.1 and Proposition~3.4]{GK})}
Let $v \in W$ and $v = s_{i_1} * \cdots * s_{i_r}$ with $i_1, \dots, i_r \in I$.
\begin{enumerate}[(a)]
\item~\label{lemma3.12_a} %[(a)]
There is a increasing sequence $1 \le k_1 < k_2 < \dots < k_l \le r$ such that 
$v = s_{i_{k_1}} s_{i_{k_2}} \dots s_{i_{k_l}}$ is a reduced expression of $v$.
In particular, we have $l(v) \le r$.
\item~\label{lemma3.12_b} %[(b)]
If $l(v) = r$, then $v = s_{i_1} \cdots s_{i_r}$.
\item~\label{lemma3.12_c} %[(c)]
If $v$ is fully commutative and $l(v) < r$, 
then there exist $a < b$ such that $s_{i_a} = s_{i_b}$ and
$s_{i_a}$ commutes with $s_{i_c}$ for every $a < c < b$.
\end{enumerate}
\end{lemm}

\begin{proof}[Proof of Proposition~\ref{prop:E*=R*}]
We denote by $\RR^*_P(F)$ the set of all subsets $D \subset P$ satisfying $w^*_D = w_F$.

First we prove $\EE^*_P(F) \subset \RR^*_P(F)$.
Since $F \in \RR^*_P(F)$, it is enough to show that, 
if $D' \in \EE^*_P(F)$ is obtained from $D \in \EE^*_P(F)$ by an ordinary or $K$-theoretical elementary excitation, 
then $w^*_{D'} = w^*_D$.
Let $u$ be a $D$-active element and $[v,u]$ be the $d_k$-interval with top element $u$.
Let $v = p_k$, $u = p_l$ and $\{ j : k < j < l, p_j \in D \} = \{ j_1, \dots, j_m \}$ ($j_1 < \dots < j_m$).
If $D' = \alpha_u(D) = D \setminus \{ u \} \cup \{ v \}$, then we have
\begin{align*}
w^*_D &= \cdots * s(p_{j_1}) * \cdots * s(p_{j_m}) * s(p_l) * \cdots,
\\
w^*_{D'} &= \cdots * s(p_k) * s(p_{j_1}) * \cdots * s(p_{j_m}) * \cdots.
\end{align*}
If $D' = \alpha^*_u(D) = D \cup \{ v \}$, then we have
\begin{align*}
w^*_D &= \cdots * s(p_{j_1}) * \cdots * s(p_{j_m}) * s(p_l) * \cdots,
\\
w^*_{D'} &= \cdots * s(p_k) * s(p_{j_1}) * \cdots * s(p_{j_m}) * s(p_l) * \cdots.
\end{align*}
Since $u$ is $D$-active, it follows from Lemma~\ref{lem:comm} and $s(p_k)*s(p_l) = s(p_l)*s(p_l) = s(p_l)$ 
that $w^*_{D'} = w^*_D$ in both cases.

Next we prove $\RR^*_P(F) \subset \EE^*_P(F)$.
We proceed by induction on $\# D$ to prove $D \in \RR^*_P(F)$ implies $D \in \EE^*_P(F)$.
Since $w^*_D = w_F$, we have $\# D \ge l(w_F) = \# F$ by Lemma~\ref{lem:red}~\ref{lemma3.12_a}.
If $\# D = \# F$, then $w^*_D = w_D$ by Lemma~\ref{lem:red}~\ref{lemma3.12_b}, 
thus we have $D \in \EE_P(F) \subset \EE^*_P(F)$ by Proposition~\ref{prop:E=R}.
If $\# D > \# F$, then it follows from Lemma~\ref{lem:red}~\ref{lemma3.12_c} that 
there exist $u$, $v \in D$ with $v < u$ such that $c(u) = c(v)$ and 
$s(u) = s(v)$ commutes with $s(w)$ for
every element $w$ between $u$ and $v$ in the expression of $w^*_D$.
If we put $D' = D \setminus \{ v \}$, then we see that $w^*_{D'} = w^*_D$.
Hence by the induction hypothesis we have $D' \in \EE^*_P(F)$.
Since $D$ is obtained from $D'$ by a sequence of ordinary and $K$-theoretical elementary excitations, 
we obtain $D \in \EE^*_P(F)$.
\end{proof}

The following proposition plays a crucial role in the proof of our main theorem 
(see the proof of Theorem~\ref{thm:xi2}).

\begin{prop}
\label{prop:E*}
Let $P$ be a connected $d$-complete poset and $F$ an order filter of $P$.
Then we have
\begin{equation}
\label{eq:E*}
\EE^*_P(F) = \bigsqcup_{D \in \EE_P(F)} \{ D \sqcup S : S \subset B(D) \}.
\end{equation}
\end{prop}

In order to prove this proposition, we prepare several lemmas.

\begin{lemm}
\label{lem:interval_BD}
Let $x$ and $y$ be elements of $P$ such that $y<x$ and $c(x) = c(y) = i$, 
and let $D\in \EE_P(F)$ be an excited diagram.
If $[y,x] \cap D \cap N_i = \emptyset$ and $y \in D$, 
then we have $[y,x] \cap B(D) \cap N_i = \emptyset$.
\end{lemm}

\begin{proof}
Assume to the contrary that $[y,x] \cap B(D) \cap N_i \neq \emptyset$ 
and take an element $z \in [y,x] \cap B(D) \cap N_i$.
By Proposition~\ref{prop:BD}~\ref{prop3.7_a}, there exists $w \in D_{c(z)}$ such that 
$w < z$ and $[w,z] \cap D \cap N_{c(z)} = \emptyset$.
Then $c(w) = c(z)$ is adjacent to $i = c(y)$ in $\Gamma$ and $w \in D \cap N_{c(x)}$.
Since $[y,x] \cap D \cap N_{c(x)} = \emptyset$, we have $w \not\in [y,x]$.
Hence by using Property~\ref{prop2.5_C5} in Proposition~\ref{prop:dc_color} 
we see that $w < y < z < x$, which contracts to $[w,z] \cap D \cap N_{c(z)} = \emptyset$.
\end{proof}

For $E \in \EE^*_P(F)$, we define a subset $S(E)$ of $E$ by putting
\begin{equation}
\label{eq:def_SE}
S(E)
 = 
\left\{
x \in E :
\begin{matrix}
\text{there exists $y \in E_{c(x)}$ such that } \\
\text{$y<x$ and $[y,x] \cap E \cap N_{c(x)} = \emptyset$}
\end{matrix}
\right\}.
\end{equation}
It follows from Property~\ref{prop2.5_C4} in Proposition~\ref{prop:dc_color} that $S(F) = \emptyset$ 
for an order filter $F$ of $P$.

\begin{lemm}
\label{lem:SE1}
Let $E \in \EE^*_P(F)$ and $u \in E$ an $E$-active element.
Then we have
\begin{align}
\label{eq:SE11}
S(\alpha_u(E))
 &=
\begin{cases}
 S(E) \setminus \{ u \} \cup \{ v \} &\text{if $u \in S(E)$,} \\
 S(E) &\text{if $u \not\in S(E)$,}
\end{cases}
\\
\label{eq:SE12}
S(\alpha^*_u(E))
 &=
\begin{cases}
 S(E) \cup \{ v \} &\text{if $u \in S(E)$,} \\
 S(E) \cup \{ u \} &\text{if $u \not\in S(E)$,}
\end{cases}
\end{align}
where $[v,u]$ is the $d_k$-interval with top element $u$.
In particular, we have
$$
\# S(\alpha_u(E)) = \# S(E),
\quad
\# S(\alpha^*_u(E)) = \# S(E) + 1,
$$
and $S(E) = \emptyset$ if and only if $E \in \EE_P(F)$.
\end{lemm}

\begin{proof}
Since $\alpha_u(E) = E \setminus \{ u \} \cup \{ v \}$, the equality~\eqref{eq:SE11} 
follows from the following three claims:
\begin{enumerate}[(i)]
\item\label{proof_lemma_3.15_i} %[(i)]
$S(\alpha_u(E)) \setminus \{ v \} \subset S(E)$.
\item\label{proof_lemma_3.15_ii} %[(ii)]
$u \in S(E)$ if and only if $v \in S(\alpha_u(E))$.
\item\label{proof_lemma_3.15_iii} %[(iii)]
$S(E) \subset S(\alpha_u(E))$.
\end{enumerate}

To prove~\ref{proof_lemma_3.15_i}, we take $x \in S(\alpha_u(E))$ such that $x \neq v$.
Then, by the definition~\eqref{eq:def_SE}, there exists $y \in (\alpha_u(E))_{c(x)}$ such that $y<x$ 
and $[y,x] \cap \alpha_u(E) \cap N_{c(x)} = \emptyset$.
If $y = v$, then $[u,x] \cap E \cap N_{c(x)} \subset [v,x] \cap E \cap N_{c(x)} = \emptyset$, 
hence $x \in S(E)$.
We consider the case where $y \neq v$.
In this case, $y \in E$ and it follows from $[v,u] \cap E \cap N_{c(u)} = \emptyset$ 
that $y \not\in [v,u] \cap N_{c(u)}$.
Now we can use Lemma~\ref{lem:interval}~\ref{lemma3.8_c} to obtain $[y,x] \cap E \cap N_{c(x)} = \emptyset$, 
hence $x \in S(E)$.

Next we prove~\ref{proof_lemma_3.15_ii}.
If $v\in S(\alpha_u(E))$, then there exists $w \in \alpha_u(E)_{c(v)}$ such that 
$w<v$ and $[w,v] \cap \alpha_u(E) \cap N_{c(v)} = \emptyset$.
Then by Lemma~\ref{lem:interval}~\ref{lemma3.8_a} we have $[w,u] \cap E \cap N_{c(u)} = \emptyset$, 
hence $u \in S(E)$.
Conversely, 
if $u \in S(E)$, there exists $z \in E_{c(u)}$ such that $z < u$ and $[z,u] \cap E \cap N_{c(u)} = \emptyset$.
Then we have $[z,v] \cap \alpha_u(E) \cap N_{c(v)} \subset [z,u] \cap \alpha_u(E) \cap N_{c(u)} = \emptyset$, 
hence $v \in S(\alpha_u(E))$.

To prove~\ref{proof_lemma_3.15_iii}, we take $x \in S(E)$.
By the definition~\eqref{eq:def_SE}, 
there exists $y \in E_{c(x)}$ such that $y<x$ and $[y,x] \cap E \cap N_{c(x)} = \emptyset$.
Then we have $y = u$ or $y \in \alpha_u(E)$.
If $y=u$, then $[u,x] \cap E \cap N_{c(x)} = \emptyset$ and 
$[v,u] \cap \alpha_u(E) \cap N_{c(u)} = \emptyset$ by the $E$-activity of $u$.
Hence by using Lemma~\ref{lem:interval}~\ref{lemma3.8_a} we have $[v,x] \cap \alpha_u(E) \cap N_{c(x)} = \emptyset$ 
and $x \in S(\alpha_u(E))$.
We consider the case where $y \in \alpha_u(E)$.
Since $[v,u] \cap E \cap N_{c(u)} = \emptyset$ and $x \in E$, we have $x \not\in [v,u] \cap N_{c(u)}$.
Hence by using Lemma~\ref{lem:interval}~\ref{lemma3.8_b} we see $[y,x] \cap \alpha_u(E) \cap N_{c(x)} = \emptyset$ 
and $x \in S(\alpha_u(E))$.

Since $\alpha^*_u(E) = E \cup \{ v \}$ and 
$u\in S(\alpha^*_u(E))$, the equality~\eqref{eq:SE12} 
follows from the following three claims:
\begin{enumerate}[(i)]\setcounter{enumi}{3}
\item\label{proof_lemma_3.15_iv} %[(iv)]
$S(\alpha^*_u(E)) \setminus \{ u, v \} \subset S(E)$.
\item\label{proof_lemma_3.15_v} %[(v)]
$u \in S(E)$ if and only if $v \in S(\alpha^*_u(E))$.
\item\label{proof_lemma_3.15_vi} %[(vi)]
$S(E) \subset S(\alpha^*_u(E))$.
\end{enumerate}

To prove~\ref{proof_lemma_3.15_iv}, we take $x \in S(\alpha^*_u(E))$ such that $x \neq u$, $v$.
Then there exists $y \in (\alpha^*_u(E))_{c(x)}$ such that $y<x$ 
and $[y,x] \cap \alpha^*_u(E) \cap N_{c(x)} = \emptyset$.
If $y = v$, then $[u,x] \cap E \cap N_{c(x)} \subset [v,x] \cap E \cap N_{c(x)} = \emptyset$, 
hence $x \in S(E)$.
If $y \in E$, then $[y,x] \cap E \cap N_{c(x)} \subset [y,x] \cap \alpha^*_u(E) \cap N_{c(x)} = \emptyset$, 
hence $x \in S(E)$.

Next we prove~\ref{proof_lemma_3.15_v}.
If $v \in S(\alpha^{*}_u(E))$, then there exists $w \in \alpha^*_u(E)_{c(v)}$ such that
$w<v$ and $[w,v] \cap \alpha^*_u(E) \cap N_{c(v)} = \emptyset$.
Then by using Lemma~\ref{lem:interval}~\ref{lemma3.8_a} we have $[w,u] \cap E \cap N_{c(u)} = \emptyset$, 
hence $u \in S(E)$.
Conversely, if $u \in S(E)$, there exists $z \in E_{c(u)}$ such that $z<u$ and 
$[z,u] \cap E \cap N_{c(u)} = \emptyset$.
Then we have $[z,v] \cap \alpha^*_u(E) \cap N_{c(v)} \subset [z,u] \cap \alpha^*_u(E) \cap N_{c(v)} = \emptyset$, 
hence $v \in S(\alpha^*_u(E))$.

To prove~\ref{proof_lemma_3.15_vi}, we take $x \in S(E)$.
Then there exists $y \in E_{c(x)}$ such that $y < x$ and $[y,x] \cap E \cap N_{c(x)} = \emptyset$.
Since $[v,u] \cap E \cap N_{c(u)} = \emptyset$ and $x \in E$, we have $x \not\in [v,u] \cap N_{c(u)}$.
Hence by using Lemma~\ref{lem:interval}~\ref{lemma3.8_b} we see $[y,x] \cap \alpha^*_u(E) \cap N_{c(x)} = \emptyset$ 
and $x \in S(\alpha^*_u(E))$.
\end{proof}

\begin{lemm}
\label{lem:SE2}
If $D \in \EE_P(F)$ and $S \subset B(D)$, 
then we have $D \sqcup S \in \EE^*_P(F)$ and $S = S(D \sqcup S)$.
In particular, if $E \in \EE^*_P(F)$ is expressed as $E = D \sqcup S = D' \sqcup S'$ 
with $D$, $D' \in \EE_P(F)$ and $S \subset B(D)$, $S' \subset B(D')$, then we have 
$D = D'$ and $S = S'$.
\end{lemm}

\begin{proof}
First we proceed by induction on $\# S$ to prove $D \sqcup S \in \EE^*_P(F)$.
If $S = \emptyset$, then we have $D \in \EE_P(F) \subset \EE^*_P(F)$.
If $S \neq \emptyset$, we take an element $x \in S$ and put $S' = S \setminus \{ x \}$.
By the induction hypothesis and Proposition~\ref{prop:E*=R*}, 
we have $D \sqcup S' \in \EE^*_P(F)$ and $w^*_{D \sqcup S'} = w_F$.
Using Proposition~\ref{prop:BD}~\ref{prop3.7_a}, we see that there exists $y \in D_{c(x)}$ 
such that $y<x$ and $[y,x] \cap D \cap N_{c(x)} = \emptyset$.
Then it follows from Lemma~\ref{lem:interval_BD} that $[y,x] \cap (D \sqcup S) \cap N_{c(x)} = \emptyset$.
If $y = p_k$, $x = p_l$ and $\{ j : k < j < l, \ p_j \in D \sqcup S \} = \{ j_1, \dots, j_m \}$ 
($j_1 < \dots < j_m$), then we have 
\begin{align*}
w^*_{D \sqcup S} &= \cdots * s(p_k) * s(p_{j_1}) * \cdots * s(p_{j_m}) * s(p_l) * \cdots,
\\
w^*_{D \sqcup S'} &= \cdots * s(p_k) * s(p_{j_1}) * \cdots * s(p_{j_m}) * \cdots.
\end{align*}
By using Lemma~\ref{lem:comm} and $s(p_k)*s(p_l) = s(p_k) * s(p_k) = s(p_k)$, 
we obtain $w^*_{D \sqcup S} = w^*_{D \sqcup S'} = w_F$.
Hence by Proposition~\ref{prop:E*=R*} we have $D \sqcup S \in \EE^*_P(F)$.

Next we put $E = D \sqcup S$ and prove that $S = S(E)$.
In order to show the inclusion $S \subset S(E)$, we take $x \in S \subset B(D)$.
Then by Proposition~\ref{prop:BD}~\ref{prop3.7_a}, there exists $y \in D_{c(x)}$ 
such that $[y,x] \cap D \cap N_{c(x)} = \emptyset$.
Hence by using Lemma~\ref{lem:interval_BD} we see that $[y,x] \cap E \cap N_{c(x)} = \emptyset$ 
and $x \in S(E)$.
In order to show the reverse inclusion $S(E) \subset S$, we take $x \in S(E)$ and prove $x \in B(D)$.
By the definition~\eqref{eq:def_SE}, there exists $y \in E_{c(x)}$ 
such that $y < x$ and $[y,x] \cap E \cap N_{c(x)} = \emptyset$.
Since $D \subset E$, we have $[y,x] \cap D \cap N_{c(x)} = \emptyset$.
If $y \in D$, then by Proposition~\ref{prop:BD}~\ref{prop3.7_a}, we have $x \in B(D)$.
If $y \in S$, then there exists $z \in D_{c(y)}$ such that $z<y$ and $[z,y] \cap D \cap N_{c(y)} = \emptyset$.
Then by using Lemma~\ref{lem:interval}~\ref{lemma3.8_a}, we have $[z,x] \cap D \cap N_{c(x)} = \emptyset$ and $x \in B(D)$.
\end{proof}

\eject
\begin{lemm}
\label{lem:SE3}
Let $E \in \EE^*_P(F)$ and $z \in S(E)$. If we put $E' = E \setminus \{ z \}$, then we have
\begin{enumerate}[(a)]
\item\label{lemma3.17_a} %[(a)]
$E' \in \EE^*_P(F)$.
\item\label{lemma3.17_b} %[(b)]
$S( E' ) = S(E) \setminus \{ z \}$.
\end{enumerate}
\end{lemm}

\begin{proof}
By the definition~\eqref{eq:def_SE}, 
there exists $w \in E_{c(z)}$ such that $w < z$ and $[w,z] \cap E \cap N_{c(z)} = \emptyset$.

%(a)
\ref{lemma3.17_a}~If $w = p_k$, $z = p_l$ and $\{ j : k < j < l, \ p_j \in E \} = \{ j_1, \dots, j_m \}$ 
($j_1 < \dots < j_m$), then we have 
\begin{align*}
w^*_E &= \cdots * s(p_k) * s(p_{j_1}) * \cdots * s(p_{j_m}) * s(p_l) * \cdots,
\\
w^*_{E'} &= \cdots * s(p_k) * s(p_{j_1}) * \cdots * s(p_{j_m}) * \cdots.
\end{align*}
By using Lemma~\ref{lem:comm} and $s(p_k)*s(p_l) = s(p_k)$, 
we obtain $w^*_E = w^*_{E'}$.
Hence it follows from Proposition~\ref{prop:E*=R*} that $E' \in \EE^*_P(F)$.

%(b)
\ref{lemma3.17_b}~First we prove that $S(E') \subset S(E)$.
Let $x \in S(E')$.
Then there exists $y \in (E')_{c(x)}$ such that $y < x $ and $[y,x] \cap E' \cap N_{c(x)} = \emptyset$.
Since $y \in E$ and $[w,z] \cap E \cap N_{c(z)} = \emptyset$, we have $y \not\in [w,z] \cap N_{c(z)}$.
Hence by using Lemma~\ref{lem:interval}~\ref{lemma3.8_d} we have $[y,x] \cap E \cap N_{c(x)} = \emptyset$ and $x \in S(E)$.

Next we prove that $S(E) \setminus \{ z \} \subset S(E')$.
We take an element $x \in S(E)$ such that $x \neq z$.
Then there exists $y \in E_{c(x)}$ such that $y<x$ and $[y,x] \cap E \cap N_{c(x)} = \emptyset$.
If $y \neq z$, then $y \in E'$ and $[y,x] \cap E' \cap N_{c(x)} = \emptyset$, hence $x \in S(E')$.
We consider the case where $y = z$.
In this case $[z,x] \cap E' \cap N_{c(x)} = \emptyset$.
Since $[w,z] \cap E' \cap N_{c(z)} = \emptyset$, it follows from Lemma~\ref{lem:interval}~\ref{lemma3.8_a} 
that $[w,x] \cap E' \cap N_{c(x)} = \emptyset$ and $x \in S(E')$.
\end{proof}

Now we are in position to give a proof of Proposition~\ref{prop:E*}.

\begin{proof}[Proof of Proposition~\ref{prop:E*}]
By using Lemma~\ref{lem:SE2}, it is enough to show that any $E \in \EE^*_P(F)$ 
can be written as $E = D \sqcup S$ with $D \in \EE_P(F)$ and $S \subset B(D)$.
Given $E \in \EE^*_P(F)$, we put $D = E \setminus S(E)$ and prove that 
$D \in \EE_P(F)$ and $S(E) \subset B(D)$.
We proceed by induction on $\# S(E)$.
If $S(E) = \emptyset$, then $E \in \EE_P(F)$ by Lemma~\ref{lem:SE1}.
We consider the case where $S(E) \neq \emptyset$.
Then we take $z \in S(E)$ and put $E' = E \setminus \{ z \}$.
Since $S(E') = S(E) \setminus \{ z \}$ by Lemma~\ref{lem:SE3}~\ref{lemma3.17_b}, 
we see that $D = E \setminus S(E) = E' \setminus S(E')$.
By the induction hypothesis, $D \in \EE_P(F)$ and $S(E') \subset B(D)$.
It remains to show that $z \in B(D)$.
By the definition~\eqref{eq:def_SE}, 
there exists $w \in E_{c(z)}$ such that $w<z$ and $[w,z] \cap E \cap N_{c(z)} = \emptyset$.
Let $w$ be the minimal such element.
If $w \not\in D$, $w \in S(E)$, 
then by definition there exists $w' \in E_{c(w)}$ such that $[w',w] \cap E \cap N_{c(w)} = \emptyset$.
Then it follows from Lemma~\ref{lem:interval}~\ref{lemma3.8_a} that $[w',z] \cap E \cap N_{c(z)} = \emptyset$, 
which contradicts to the minimality of $w$.
Therefore we have $w \in D$ and $z \in B(D)$.
\end{proof}

\begin{rema}
\label{rem:heap2}
We can define the notion of $D$-active elements and ordinary and $K$-theoretical elementary excitations 
for a general dominant minuscule heap $H(w)$ just by replacing $d_k$-intervals with 
intervals $[v,u]$ such that $c(u) = c(v)$ and $[v,u] \cap H(w)_{c(u)} = \{ u, v \}$.
Then the arguments in this section work as well for $H(w)$.
In particular, Propositions~\ref{prop:E*=R*} and~\ref{prop:E*} holds for $H(w)$.
\end{rema}

\begin{exam}
\label{ex:shited-B-excited}
Let $\mu$ be a strict partition and regard the shifted shape $S(\mu)$ as a heap for the Weyl group of type $B$ 
(see Example~\ref{ex:shifted-B}).
Then the notion of active elements, (ordinary and $K$-theoretical) elementary excitations and excited peaks 
are modified as follows.
Let $D$ be a subset of $S(\mu)$.
\begin{enumerate}[(a)]
\item\label{exam3.19_a} %[(a)]
An element $u = (i,j) \in D$ is $D$-active if either
\begin{gather*}
\text{$i < j$ and $(i,j+1)$, $(i+1,j)$, $(i+1,j+1) \in S(\mu) \setminus D$, or}
\\